% ================================================================
% PROCESSI DI SUPPORTO
% ================================================================

\section{Processi Di Supporto}
\label{sec:processi-supporto}

In conformità allo standard ISO/IEC~12207:1995\textsubscript{G}, i processi di
supporto comprendono le attività trasversali necessarie a garantire qualità,
tracciabilità, coerenza e controllo degli artefatti prodotti.

Nel contesto del progetto \textit{Second Brain}, e in accordo con il tailoring
descritto nella Sezione~\ref{sec:tailoring}, il gruppo \textit{Synergo Unit} adotta
i seguenti processi di supporto:

\begin{itemize}
    \item documentazione;
    \item gestione della configurazione;
    \item verifica;
    \item validazione;
    \item accertamento della qualità;
    \item gestione dei cambiamenti;
    \item gestione delle non conformità;
\end{itemize}

% ================================================================
% DOCUMENTAZIONE
% ================================================================

\subsection{Processo Di Documentazione}
\label{sec:documentazione}

\subsubsection{Descrizione}

Il processo di documentazione definisce regole, strumenti e ciclo di vita per la
produzione, la revisione, l'approvazione e la pubblicazione della documentazione
del gruppo.

Esso si applica a ogni documento gestionale o tecnico, interno o esterno, prodotto
dal gruppo \textit{Synergo Unit}. Il Diario di Bordo segue un template separato
realizzato tramite Canva\textsubscript{G} e non prevede registro delle modifiche.

La documentazione deve essere:
\begin{itemize}
    \item uniforme nella struttura;
    \item coerente con le presenti Norme di Progetto;
    \item verificabile tramite Pull Request\textsubscript{G};
    \item tracciabile rispetto alle modifiche logiche effettuate;
    \item consultabile tramite repository\textsubscript{G} e sito documentale.
\end{itemize}

% ------------------------------------------------

\subsubsection{Tipologia Dei Documenti}

\paragraph{Documentazione Gestionale Interna}

\begin{itemize}
    \item \textbf{Norme di Progetto}: definiscono processi, convenzioni, strumenti
    e procedure operative adottate dal gruppo. Sono redatte e mantenute
    dall'Amministratore\textsubscript{G}.

    \item \textbf{Glossario}: raccoglie termini tecnici, acronimi ed espressioni
    ambigue utilizzate nella documentazione. Ogni autore può proporre nuovi
    termini; l'Amministratore ne cura normalizzazione e coerenza.

    \item \textbf{Analisi dello Stato dell'Arte}: documento che raccoglie valutazione comparativa 
    delle tecnologie e motivazione dello stack candidato per il Proof of Concept\textsubscript{G}.

    \item \textbf{Verbali Interni}: documentano riunioni interne, decisioni operative,
    retrospettive, pianificazione e assegnazione dei task\textsubscript{G}.
\end{itemize}

\paragraph{Documentazione Gestionale Esterna}

\begin{itemize}
    \item \textbf{Verbali Esterni}: documentano gli incontri con la proponente e con
    gli stakeholder\textsubscript{G} esterni. Sono redatti dall'Amministratore e
    condivisi con la proponente per conferma e firma.

    \item \textbf{Piano di Progetto}: descrive pianificazione, consuntivazione,
    gestione dei rischi, controllo dei costi e rotazione dei ruoli.

    \item \textbf{Piano di Qualifica}: definisce strategie, metriche e attività di
    verifica e validazione adottate dal gruppo.

    \item \textbf{Diario di Bordo}: materiale gestionale usato per il confronto
    periodico con il committente. È redatto tramite template Canva e non segue
    il ciclo di versionamento documentale dei documenti \LaTeX{}.

    \item \textbf{Sito Web}: punto di accesso alla documentazione pubblicata tramite
    GitHub Pages\textsubscript{G}.
\end{itemize}

\paragraph{Documentazione Tecnica}

\begin{itemize}
    \item \textbf{Analisi dei Requisiti}: documento tecnico esterno che formalizza
    attori, casi d'uso, requisiti, fonti e matrici di tracciabilità.
\end{itemize}

% ------------------------------------------------

\subsubsection{Strumenti Di Documentazione}
\label{sec:strumenti-doc}

La documentazione viene prodotta tramite:

\begin{itemize}
    \item \textbf{\LaTeX{} con Visual Studio Code}: ambiente principale di redazione;

    \item \textbf{LaTeX Workshop}: estensione utilizzata per compilare localmente i
    documenti PDF;

    \item \textbf{Repository GitHub}: archivio ufficiale della documentazione e della
    relativa storia;

    \item \textbf{GitHub Pages}: sito documentale aggiornato periodicamente tramite
    task dedicato;

    \item \textbf{Canva}: strumento utilizzato per il template dei Diari di Bordo;

    \item \textbf{Script Python per Glossario}: strumento utilizzato per applicare
    automaticamente il pedice \textsubscript{G} alla prima occorrenza dei termini
    presenti nel Glossario.
\end{itemize}

Ogni membro lavora esclusivamente su copia locale del repository.

% ------------------------------------------------

\subsubsection{Convenzioni Redazionali}
\label{sec:convenzioni}

Al fine di garantire uniformità e coerenza tra i documenti, il gruppo adotta le
seguenti convenzioni:

\begin{itemize}
    \item \textbf{Nomi dei file}: tutti i nomi devono essere scritti in minuscolo,
    utilizzando il carattere \texttt{\_} come separatore di parola.
    \\
    \textit{Esempio}: \texttt{piano\_di\_progetto.tex}.

    \item \textbf{Nome del file principale}: il file principale \LaTeX{} deve essere
    nominato secondo il nome previsto per il PDF finale.
    \\
    \textit{Esempio}: \texttt{norme\_di\_progetto.tex} produce
    \texttt{norme\_di\_progetto.pdf}.

    \item \textbf{Titoli}: sezioni, sottosezioni e paragrafi devono utilizzare
    Title Case\textsubscript{G}.

    \item \textbf{Tabelle lunghe}: le tabelle che possono superare una pagina devono
    essere realizzate tramite \texttt{xltabular} o \texttt{longtable}, in modo da
    consentirne la spezzatura automatica.

    \item \textbf{Pedice G}: la prima occorrenza, nell'intero documento, di ogni
    termine presente nel Glossario deve essere marcata con il pedice
    \textsubscript{G}. Lo script riceve in ingresso la concatenazione dei file
    \LaTeX{} delle sezioni e applica il pedice alla prima occorrenza globale.

    \item \textbf{Link generici}: i collegamenti ipertestuali nel testo devono essere
    prodotti tramite il comando \verb|\link{url}{testo}|.

    \item \textbf{Link documentali in tabelle}: i collegamenti a documenti all'interno
    di tabelle devono essere prodotti tramite il comando
    \verb|\doclink{url}{testo}|, così da preservare leggibilità e andata a capo.

    \item \textbf{Issue GitHub}: i riferimenti a issue devono essere prodotti tramite
    il comando \verb|\ghissue{id}|.

    \item \textbf{Decisioni da verbale interno}: le decisioni interne devono essere
    identificate con codice \texttt{VI-y.x}, dove \texttt{y} è il numero progressivo
    del verbale e \texttt{x} è il numero progressivo della decisione.

    \item \textbf{Decisioni da verbale esterno}: le decisioni esterne devono essere
    identificate con codice \texttt{VE-y.x}, con la stessa logica adottata per i
    verbali interni.
\end{itemize}

% ------------------------------------------------

\subsubsection{Template Documentale}

Tutti i documenti, ad eccezione dei Diari di Bordo, condividono un template
\LaTeX{} comune comprensivo di:

\begin{itemize}
    \item frontespizio;
    \item registro delle modifiche;
    \item indice;
    \item sezioni specifiche del documento;
    \item configurazioni condivise;
    \item loghi e asset comuni.
\end{itemize}

Il template definisce inoltre:
\begin{itemize}
    \item margini;
    \item colori;
    \item profondità dell'indice;
    \item comandi personalizzati;
    \item configurazione dei link;
    \item formattazione delle tabelle;
    \item stile dei paragrafi.
\end{itemize}

% ------------------------------------------------

\subsubsection{Registro Delle Modifiche}

Ogni documento \LaTeX{}, ad eccezione dei Diari di Bordo, deve contenere un registro
delle modifiche.

Il registro deve riportare:

\begin{itemize}
    \item \textbf{Versione};
    \item \textbf{Data};
    \item \textbf{Autore};
    \item \textbf{Verificatore};
    \item \textbf{Descrizione}.
\end{itemize}

Il registro delle modifiche è aggiornato dal Redattore\textsubscript{G}.

Le voci del registro devono descrivere modifiche logiche e documentali significative.

% ------------------------------------------------

\subsubsection{Gestione Del Glossario}

Il Glossario viene aggiornato in modo incrementale durante la produzione della
documentazione.

Il processo prevede:

\begin{enumerate}
    \item durante la stesura, ogni autore individua eventuali nuovi termini tecnici,
    acronimi o espressioni ambigue;

    \item al termine della stesura, l'autore inserisce i nuovi termini nel file
    \texttt{termini.tex} del Glossario;

    \item l'autore aggiorna il file \texttt{glossary.txt}, utilizzato dallo script
    Python;

    \item l'autore esegue manualmente lo script per applicare il pedice
    \textsubscript{G} alla prima occorrenza dei termini nel documento;

    \item il Verificatore controlla che i termini inseriti siano coerenti con il
    documento e che il pedice sia stato applicato correttamente.
\end{enumerate}

L'Amministratore può intervenire per normalizzare definizioni, duplicati,
maiuscole/minuscole e sinonimi.

% ------------------------------------------------

\subsubsection{Checklist Prima Della Pull Request}

Prima di aprire una Pull Request, il Redattore deve verificare che:

\begin{itemize}
    \item il documento compili senza errori bloccanti;
    \item il PDF sia generato correttamente;
    \item il registro delle modifiche sia aggiornato, così come numero di versione e data di ultima modifica;
    \item non siano presenti segnaposto non motivati;
    \item i link siano funzionanti;
    \item i riferimenti interni siano corretti;
    \item le tabelle lunghe siano spezzabili;
    \item il Glossario sia aggiornato, se necessario;
    \item il pedice \textsubscript{G} sia stato applicato correttamente;
    \item il nome del file sia conforme alle convenzioni.
\end{itemize}

% ------------------------------------------------

\subsubsection{Ciclo Di Vita Documentale}
\label{sec:ciclo-vita}

Ogni documento attraversa i seguenti stati:

\begin{enumerate}
    \item \textbf{In Stesura}: il Redattore produce o modifica il contenuto su branch
    dedicato;

    \item \textbf{In Verifica}: il Redattore apre una Pull Request; il Verificatore
    controlla il documento e può approvare o richiedere modifiche;

    \item \textbf{Verificato}: il documento ha superato la verifica tramite Pull Request;

    \item \textbf{Approvato}: il documento verificato viene approvato dal Responsabile,
    eventualmente dopo review di gruppo;

    \item \textbf{Rilasciato}: il documento approvato viene pubblicato nel repository
    e reso disponibile tramite sito documentale.
\end{enumerate}

Un documento entra in baseline solo dopo approvazione del Responsabile e review
di gruppo.

% ------------------------------------------------

\subsubsection{Versionamento Documentale}

Il gruppo adotta il \textit{Semantic Versioning}\textsubscript{G} a tre cifre nel
formato:

\begin{center}
    \texttt{vX.Y.Z}
\end{center}

Le versioni vengono incrementate secondo le seguenti regole:

\begin{xltabular}{\textwidth}{
    >{\raggedright\arraybackslash}p{2.5cm}
    >{\raggedright\arraybackslash}p{3cm}
    >{\raggedright\arraybackslash}X
}
    \toprule
    \textbf{Campo} & \textbf{Incremento} & \textbf{Quando Applicarlo} \\
    \midrule
    \endfirsthead

    \toprule
    \textbf{Campo} & \textbf{Incremento} & \textbf{Quando Applicarlo} \\
    \midrule
    \endhead

    \midrule
    \multicolumn{3}{r}{\textit{Continua nella pagina successiva}} \\
    \endfoot

    \bottomrule
    \endlastfoot

    Major
    & \texttt{+1.0.0}
    & Ingresso in una nuova baseline o cambiamento strutturale profondo del documento. \\

    Minor
    & \texttt{+0.1.0}
    & Aggiunta o modifica significativa di contenuti. \\

    Patch
    & \texttt{+0.0.1}
    & Correzioni ortografiche, tipografiche o modifiche non semantiche. \\

\end{xltabular}

L'incremento di versione è scelto dal Redattore, controllato dal Verificatore e
approvato dal Responsabile in caso di rilascio ufficiale.

% ================================================================
% GESTIONE DELLA CONFIGURAZIONE
% ================================================================

\subsection{Processo Di Gestione Della Configurazione}
\label{sec:configurazione}

\subsubsection{Descrizione}

Il processo di gestione della configurazione garantisce integrità, consistenza,
versionamento e tracciabilità degli elementi configurativi del progetto.

Sono considerati elementi configurativi:

\begin{itemize}
    \item documenti \LaTeX{};
    \item file PDF generati;
    \item template;
    \item asset grafici;
    \item script di automazione;
    \item codice del Proof of Concept, quando presente;
    \item configurazioni del repository;
    \item materiale pubblicato su GitHub Pages.
\end{itemize}

% ------------------------------------------------

\subsubsection{Sistema Di Gestione}

Il gruppo utilizza GitHub come sistema principale di gestione della configurazione.

Ogni modifica deve essere tracciata tramite:

\begin{itemize}
    \item issue;
    \item branch dedicato;
    \item commit riferiti alla issue;
    \item Pull Request;
    \item review;
    \item merge controllato.
\end{itemize}

% ------------------------------------------------

\subsubsection{Struttura Del Repository}

La repository è organizzata per baseline e tipologia documentale.

La struttura deve mantenere separati:

\begin{itemize}
    \item documenti interni;
    \item documenti esterni;
    \item verbali interni;
    \item verbali esterni;
    \item documentazione tecnica;
    \item template;
    \item asset condivisi;
    \item automazioni.
\end{itemize}

La versione ufficiale e stabile di ogni elemento configurativo è quella presente
nel branch \texttt{main}.

% ------------------------------------------------

\subsubsection{Strategia Di Branching}

Il gruppo adotta un modello \textit{trunk-based development}\textsubscript{G}.

Il branch \texttt{main} è l'unico branch stabile. Ogni modifica deve essere sviluppata su branch dedicato, creato a partire dalla
sezione \textit{Development} della issue GitHub associata.

La convenzione di naming dei branch è:

\begin{center}
    \texttt{<id>-<[gestionale|tecnico]>\_<[esterno|interno]>\_<nomeDoc>}
\end{center}

\textit{Esempio}: \texttt{116-gestionale\_interno\_norme\_di\_progetto}

% ------------------------------------------------

\subsubsection{Gestione Delle Issue}

Le issue devono essere create esclusivamente tramite Google Forms.

Non è consentita la creazione manuale di issue, salvo futura modifica esplicita
delle presenti norme.

Ogni issue deve contenere:

\begin{itemize}
    \item assegnatario;
    \item ruolo;
    \item priorità;
    \item scadenza;
    \item tempo previsto;
    \item milestone RTB o PB;
    \item descrizione dell'attività.
\end{itemize}

% ------------------------------------------------

\subsubsection{Priorità Delle Issue}

Le issue vengono classificate secondo:

\begin{itemize}
    \item \textbf{Bassa}: attività non urgente o non vincolante per lo sprint corrente;
    \item \textbf{Media}: priorità standard;
    \item \textbf{Alta}: attività da completare entro lo sprint corrente;
    \item \textbf{Urgente}: criticità che richiede intervento immediato.
\end{itemize}

% ------------------------------------------------

\subsubsection{Workflow Delle Issue}

Le issue attraversano i seguenti stati all'interno di GitHub Projects:

\begin{itemize}
    \item \textbf{To Do};
    \item \textbf{In Progress};
    \item \textbf{In Review};
    \item \textbf{Done}.
\end{itemize}

Le transizioni sono gestite tramite automazioni:

\begin{itemize}
      \item la creazione del branch dalla issue porta la issue in \textit{In Progress};
      \item l'apertura della Pull Request porta la issue in \textit{In Review};
      \item il merge della Pull Request chiude automaticamente la issue e la porta in \textit{Done}.
\end{itemize}

% ------------------------------------------------

\subsubsection{Commit}

Ogni commit relativo a una issue deve contenere il riferimento all'identificativo
della issue nel messaggio.

La forma prescritta è:

\begin{center}
    \texttt{\#id messaggio}
\end{center}

\textit{Esempio}: \texttt{\#116 aggiornata sezione documentazione}

% ------------------------------------------------

\subsubsection{Pull Request}

Ogni modifica deve essere integrata tramite Pull Request.

La Pull Request deve:

\begin{itemize}
      \item essere collegata alla issue corrispondente;
      \item descrivere sinteticamente le modifiche effettuate;
      \item essere aperta, salvo urgenze motivate, entro la domenica precedente la riunione del martedì;    \item ricevere almeno una review positiva;
      \item non essere approvata dall'autore della modifica.
\end{itemize}

In caso di modifiche sostanziali richieste dal Verificatore, il Redattore può
segnare la Pull Request come \textit{draft} fino al completamento delle correzioni.

% ------------------------------------------------

\subsubsection{Branch Protection Rules}

Il branch \texttt{main} è protetto tramite branch protection rules.

Le regole attive prevedono:

\begin{itemize}
    \item obbligo di almeno una review positiva;
    \item blocco dei force push;
    \item restrizione dell'eliminazione del branch protetto.
\end{itemize}

% ------------------------------------------------

\subsubsection{Merge}

Il merge su \texttt{main} può essere eseguito solo dopo approvazione della Pull
Request da parte del Verificatore.

Il merge viene eseguito dal Redattore, non dal Verificatore.

Il gruppo utilizza la modalità standard di merge proposta da GitHub tramite Pull
Request.

Dopo il merge:

\begin{itemize}
    \item la issue associata viene chiusa automaticamente;
    \item il branch di lavoro viene eliminato;
    \item il Redattore aggiorna il tempo effettivo nel foglio di consuntivazione.
\end{itemize}

% ------------------------------------------------

\subsubsection{Generazione E Pubblicazione Dei PDF}

I PDF vengono generati localmente tramite LaTeX Workshop.

Il sito GitHub Pages non viene aggiornato automaticamente a ogni merge su
\texttt{main}; l'aggiornamento del sito documentale viene gestito tramite task
periodica dedicata.

% ================================================================
% VERIFICA
% ================================================================

\subsection{Processo Di Verifica}
\label{sec:verifica}

\subsubsection{Descrizione}

Il processo di verifica assicura che ogni artefatto prodotto sia conforme alle
presenti norme e agli obiettivi di qualità definiti nel Piano di Qualifica.

La verifica risponde alla domanda:

\begin{center}
    \textit{Il prodotto o artefatto è stato realizzato correttamente rispetto alle norme stabilite?}
\end{center}

In fase RTB la verifica riguarda principalmente:

\begin{itemize}
    \item documentazione;
    \item coerenza dei riferimenti;
    \item correttezza delle matrici di tracciabilità;
    \item rispetto delle convenzioni;
    \item adeguatezza del materiale relativo al Proof of Concept.
\end{itemize}

% ------------------------------------------------

\subsubsection{Responsabilità}

\begin{itemize}
    \item Il Verificatore deve essere sempre un membro diverso dal Redattore.
    \item Nessun membro può verificare e approvare il proprio lavoro.
    \item Il Verificatore non modifica direttamente il documento verificato.
    \item Il Verificatore segnala le modifiche richieste tramite commenti o richieste
    di cambiamento nella Pull Request.
    \item Il Responsabile interviene nell'approvazione finale, non nella verifica
    puntuale del documento.
\end{itemize}

% ------------------------------------------------

\subsubsection{Attività Di Verifica}

Il Verificatore deve controllare:

\begin{itemize}
    \item correttezza linguistica;
    \item coerenza contenutistica;
    \item rispetto del template;
    \item correttezza del registro delle modifiche;
    \item correttezza dei riferimenti interni;
    \item funzionamento dei link;
    \item presenza del pedice \textsubscript{G} dove necessario;
    \item coerenza con il Glossario;
    \item conformità alle convenzioni di naming;
    \item correttezza del versionamento;
    \item compilazione corretta del documento PDF;
    \item assenza di segnaposto non motivati.
\end{itemize}

Se la modifica coinvolge documenti collegati, il Verificatore deve controllare
anche i riferimenti coinvolti.

In particolare, modifiche ai requisiti devono comportare il controllo di:

\begin{itemize}
    \item Analisi dei Requisiti;
    \item matrici di tracciabilità;
    \item Glossario, se necessario;
    \item eventuali riferimenti nel Piano di Qualifica.
\end{itemize}

% ------------------------------------------------

\subsubsection{Esito Della Verifica}

La verifica può avere:

\begin{itemize}
    \item \textbf{Esito positivo}: il Verificatore approva la Pull Request;
    \item \textbf{Esito negativo}: il Verificatore richiede modifiche tramite commenti
    o richiesta di cambiamento.
\end{itemize}

Una richiesta di modifica può bloccare il merge finché il Redattore non ha risolto
le osservazioni.

% ------------------------------------------------

\subsubsection{Definition Of Done}
\label{sec:dod}

Un task è considerato \textit{Done} solo se:

\begin{itemize}
    \item è stato creato tramite Google Form;
    \item è associato a una issue GitHub;
    \item è collegato, se necessario, a una decisione verbalizzata;
    \item è stato sviluppato su branch dedicato;
    \item i commit riportano il riferimento alla issue;
    \item la Pull Request è stata aperta, salvo urgenze motivate, entro la domenica precedente la riunione;
    \item il contenuto è completo rispetto allo scope della issue;
    \item il registro delle modifiche è stato aggiornato;
    \item il documento compila correttamente;
    \item il Verificatore ha approvato la Pull Request;
    \item il Redattore ha eseguito il merge;
    \item la issue è stata chiusa;
    \item il tempo effettivo è stato registrato.
\end{itemize}

% ================================================================
% VALIDAZIONE
% ================================================================

\subsection{Processo Di Validazione}
\label{sec:validazione}

\subsubsection{Descrizione}

Il processo di validazione ha lo scopo di accertare che quanto prodotto risponda
alle esigenze della proponente e agli obiettivi del capitolato.

La validazione risponde alla domanda:

\begin{center}
    \textit{Stiamo costruendo il prodotto giusto?}
\end{center}

In fase RTB, la validazione riguarda principalmente:

\begin{itemize}
    \item requisiti;
    \item casi d'uso;
    \item scelte tecnologiche;
    \item perimetro del Proof of Concept.
\end{itemize}

% ------------------------------------------------

\subsubsection{Validazione Con La Proponente}

Ogni incontro di validazione con la proponente deve produrre un verbale esterno.

La issue relativa alla redazione del verbale esterno rimane aperta fino a quando
il verbale non è stato firmato o confermato dalla proponente via email.

Se la proponente richiede modifiche rilevanti, queste devono essere:

\begin{itemize}
    \item riportate nel verbale esterno;
    \item tradotte in issue dedicate;
    \item implementate tramite Pull Request;
    \item verificate dal Verificatore;
    \item integrate nei documenti coinvolti.
\end{itemize}

Le modifiche ai requisiti derivanti da feedback esterno devono aggiornare:

\begin{itemize}
    \item Analisi dei Requisiti;
    \item matrici di tracciabilità;
    \item Glossario, se necessario;
    \item Piano di Qualifica, se impattano test, metriche o criteri di accettazione.
\end{itemize}

% ================================================================
% ACCERTAMENTO DELLA QUALITÀ
% ================================================================

\subsection{Processo Di Accertamento Della Qualità}
\label{sec:qualita}

\subsubsection{Descrizione}

Il processo di accertamento della qualità ha lo scopo di monitorare e migliorare
la qualità dei processi e degli artefatti prodotti.

A differenza della verifica, che controlla puntualmente un artefatto, l'accertamento
della qualità monitora l'efficacia e l'efficienza del way of working\textsubscript{G}.

% ------------------------------------------------

\subsubsection{Modello Di Miglioramento}

Il gruppo adotta il ciclo PDCA\textsubscript{G} come riferimento per il miglioramento
continuo:

\begin{itemize}
    \item \textbf{Plan}: definizione di obiettivi, metriche e procedure;
    \item \textbf{Do}: applicazione dei processi definiti;
    \item \textbf{Check}: misurazione dei risultati e analisi degli scostamenti;
    \item \textbf{Act}: introduzione di azioni correttive e aggiornamento del way of working.
\end{itemize}

% ------------------------------------------------

\subsubsection{Metriche E Piano Di Qualifica}

Le metriche, le soglie e il cruscotto di qualità sono definiti progressivamente
nel Piano di Qualifica.

Poiché in fase RTB il Piano di Qualifica è ancora in stesura, le presenti norme
definiscono il processo generale, demandando al Piano di Qualifica la definizione
di dettaglio di:

\begin{itemize}
    \item metriche di qualità di processo;
    \item metriche di qualità di prodotto;
    \item soglie di accettabilità;
    \item strumenti di misurazione;
    \item modalità di raccolta dei dati.
\end{itemize}

Ogni modifica al way of working emersa da retrospettive o azioni correttive deve
comportare, se necessario, l'aggiornamento delle Norme di Progetto.

% ================================================================
% GESTIONE DEI CAMBIAMENTI
% ================================================================

\subsection{Processo Di Gestione Dei Cambiamenti}
\label{sec:cambiamenti}

\subsubsection{Descrizione}

Il processo di gestione dei cambiamenti disciplina la valutazione, l'approvazione,
l'implementazione e la verifica delle modifiche rilevanti a documenti, requisiti,
tecnologie e processi.

Ogni cambiamento rilevante deve essere tracciato.

% ------------------------------------------------

\subsubsection{Tipologie Di Cambiamento}

I cambiamenti vengono distinti in:

\begin{itemize}
    \item \textbf{documentali}: modifiche a struttura o contenuto dei documenti;
    \item \textbf{di requisito}: modifiche a requisiti, casi d'uso o matrici di tracciabilità;
    \item \textbf{tecnologici}: modifiche allo stack o alle tecnologie candidate;
    \item \textbf{di processo}: modifiche al way of working o alle presenti norme.
\end{itemize}

% ------------------------------------------------

\subsubsection{Procedura}

Ogni cambiamento segue la seguente procedura:

\begin{enumerate}
    \item \textbf{Richiesta}: il cambiamento emerge da verbale interno, verbale esterno,
      retrospettiva, feedback della proponente o revisione esterna;

    \item \textbf{Valutazione Di Impatto}: il gruppo valuta impatto su documenti,
    requisiti, tempi, costi, tecnologie o processi;

    \item \textbf{Approvazione}: il cambiamento viene approvato dal ruolo competente;

    \item \textbf{Implementazione}: viene creata una issue e il cambiamento viene
    implementato tramite branch e Pull Request;

    \item \textbf{Verifica}: il Verificatore controlla la conformità della modifica;

    \item \textbf{Chiusura}: la issue viene chiusa dopo merge e aggiornamento del
    tempo effettivo.
\end{enumerate}

% ------------------------------------------------

\subsubsection{Responsabilità Di Approvazione}

\begin{itemize}
    \item I cambiamenti documentali ordinari sono approvati tramite verifica della
    Pull Request.

    \item I cambiamenti ai requisiti sono approvati dall'Analista e dal Responsabile;
    se incidono sul perimetro funzionale, devono essere validati con la proponente.

    \item I cambiamenti tecnologici sono approvati dal Progettista e dal Responsabile;
    se rilevanti, devono essere discussi con la proponente.

    \item I cambiamenti alle Norme di Progetto sono approvati dall'Amministratore e
    dal Responsabile.
\end{itemize}

% ================================================================
% GESTIONE DELLE NON CONFORMITÀ
% ================================================================

\subsection{Processo Di Gestione Delle Non Conformità}
\label{sec:non-conformita}

\subsubsection{Descrizione}

Una non conformità è una deviazione rispetto alle presenti norme, al template, al
Piano di Qualifica o agli accordi formalizzati nei verbali.

Le non conformità possono emergere durante:

\begin{itemize}
    \item verifica tramite Pull Request;
    \item review di sprint;
    \item retrospettiva;
    \item revisione con committente;
    \item incontro con la proponente;
    \item controllo di un documento già rilasciato.
\end{itemize}

% ------------------------------------------------

\subsubsection{Classificazione}

Le non conformità vengono classificate in:

\begin{itemize}
    \item \textbf{Bloccante}: impedisce il merge, il rilascio o l'ingresso in baseline
    dell'artefatto;

    \item \textbf{Non Bloccante}: non impedisce il merge immediato, ma deve essere
    corretta tramite task tracciato nello sprint corrente o successivo.
\end{itemize}

% ------------------------------------------------

\subsubsection{Gestione}

Se la non conformità viene individuata prima del merge, il Verificatore la segnala
nella Pull Request tramite commento o richiesta di cambiamento.

Se la non conformità viene individuata dopo il rilascio, essa deve essere:

\begin{itemize}
    \item verbalizzata nella riunione successiva;
    \item trasformata in task;
    \item corretta tramite issue dedicata;
    \item verificata tramite Pull Request;
    \item registrata nel changelog del documento, se comporta modifica documentale.
\end{itemize}

Una non conformità bloccante può impedire il rilascio di una baseline fino alla
sua risoluzione.
